
We present evidence through three analyses: (1) implementation quality versus execution feasibility comparison, (2) memory requirement breakdown demonstrating estimation complexity, and (3) timeline cost analysis.

\subsubsection{Implementation Quality vs Execution Feasibility}

Table 1 presents the central observation: all traditional quality metrics passed while the experiment was completely unrunnable.

\textbf{Table 1: Implementation Quality Metrics vs Execution Feasibility}

\begin{table}[htbp]
\centering
\begin{tabular}{llll}
\toprule
Dimension & Metric & Status & Evidence \\
\midrule
**Implementation Quality** &  &  &  \\
Task Completion & 10/10 tasks & Complete & All Phase 4 tasks finished \\
SDD Compliance & 100\% & Pass & All tasks passed TEST→IMPL→VERIFY cycle \\
Code Artifacts & 39 files & Complete & 29 Python + 10 test files \\
Lines of Code & ~8,200 LOC & Complete & Full implementation \\
Test Coverage & 10 test suites & All Passing & Unit tests for all components \\
Code Quality & High & Pass & Modular design, error handling, type hints \\
**Execution Feasibility** &  &  &  \\
Experiment Runs & 0 & Fail & Could not execute a single training step \\
Model Loading & Failed & Fail & CUDA Out of Memory during initialization \\
Memory Availability & Insufficient & Fail & Required ~489GB, available 475GB \\
Validation Gate & Failed & Fail & Hypothesis validation failed due to infeasibility \\
\bottomrule
\end{tabular}
\end{table}

Traditional software quality metrics (task completion, test passing, SDD compliance for implementation phases) do not detect computational feasibility issues. The failure occurred after all quality gates passed, indicating that feasibility is orthogonal to implementation quality.

\subsubsection{Memory Requirement Analysis}

The substantial difference between naive and realistic calculations demonstrates why systematic checking may be necessary:

```
Naive Calculation (INCORRECT):
├─ Model Parameters:        94 GB
└─ Optimizer States:       188 GB
   ────────────────────────────
   Total (naive):          282 GB  ← Appeared feasible for 475GB capacity

Realistic Calculation (CORRECT):
├─ Model Parameters:        94 GB
├─ Optimizer States:       188 GB
├─ Gradients:               94 GB
├─ Activations:             75 GB
└─ Framework Overhead:      38 GB
   ────────────────────────────
   Total (realistic):      489 GB  ← Exceeds 475GB capacity

Underestimation: 207 GB (73\% of naive estimate)
Naive captures only 58\% of actual requirement (282/489)
```

\textbf{Component distribution:} Optimizer states dominate (38\% of total), followed by model/gradients (19\% each), activations (15\%), and framework overhead (8\%). The optimizer multiplier (2× for AdamW) is the primary driver of underestimation.

\textbf{Validation:} Post-hoc analysis of CUDA memory error logs confirmed actual requirement ~490GB, matching the realistic estimate within 1GB (0.2\% error). This validates the estimation formula accuracy in this case.

\subsubsection{Timeline Cost Analysis}

Figure 1 compares actual timeline (no feasibility gate) with proposed timeline (with gate).

\textbf{Actual Timeline (No Feasibility Gate):}

```
Phase 2C: Experiment Design
│ Duration: 1-2 hours
│ Output: Mixtral-8x7B specification
└─> Approved (no feasibility check)

Phase 3: Implementation Planning  
│ Duration: 2-3 hours
│ Output: PRD, Architecture, Logic, Config
└─> Complete

Phase 4: Coding \& Testing
│ Duration: 8-13 hours  
│ Output: 29 Python files, 10 tests, 8200 LOC
└─> Complete (100\% task completion)

Phase 4 End: Execution Attempted
│ Duration: Immediate
└─> INFEASIBILITY DISCOVERED (CUDA OOM)

Total time before discovering infeasibility: 10-16 hours
```

\textbf{Proposed Timeline (With Feasibility Gate):}

```
Phase 2C: Experiment Design
│ Duration: 1-2 hours
│ Output: Mixtral-8x7B specification
└─> Pending feasibility check

Phase 2C.5: Feasibility Gate (NEW)
│ Duration: ~5 minutes
│ Actions: 
│   - Calculate: 94+188+94+75+38 = 489GB
│   - Compare: 489GB vs 475GB available
│   - Status: FAIL (103\% utilization)
└─> INFEASIBILITY DETECTED

Gate Output: Block Phase 3, surface options
│ Options:
│   1. Scale-down model (Phi-2, GPT-2)
│   2. Enable 8-bit quantization
│   3. Justify additional hardware
└─> Reformulation before implementation

Potential time saved: 10-16 hours (in this case)
Gate cost: 5 minutes (<1\% overhead)
```

\textbf{Cost-benefit ratio:} 120:1 to 192:1 (assuming gate prevents all downstream waste and no false positive costs).

\textbf{Caveat:} This ratio assumes perfect compliance (researcher heeds warning and reformulates) and negligible false positive costs (gate does not reject feasible experiments unnecessarily). Actual benefit depends on gate accuracy and researcher behavior.

\subsubsection{Retrospective Gate Performance}

Table 2 evaluates gate accuracy on the failure case and hypothetical alternatives.

\textbf{Table 2: Gate Accuracy on Configurations}

\begin{table}[htbp]
\centering
\begin{tabular}{lllll}
\toprule
Configuration & True Status & Gate Prediction & Memory Est. & Accuracy \\
\midrule
Mixtral-8x7B (actual) & Infeasible & Infeasible ✓ & 489 GB & Correct (0.2\% error vs ~490GB actual) \\
Phi-2 (2.7B params) & Feasible* & Feasible ✓ & 25 GB & Correct* \\
GPT-2 XL (1.5B) & Feasible* & Feasible ✓ & 15 GB & Correct* \\
Mixtral + 8-bit quant & Feasible* & Feasible ✓ & 350 GB & Correct* \\
5 tasks (reduced) & Infeasible* & Infeasible ✓ & 455 GB & Correct* \\
\bottomrule
\end{tabular}
\end{table}

*Hypothetical - not empirically tested

The gate correctly identified the actual infeasible configuration with high estimation accuracy (489GB estimate vs ~490GB actual requirement = 0.2\% error). Hypothetical alternatives were not empirically verified. Statistical validation requires testing across diverse configurations.
